\appendix
\section{An independent split-basis cubic calculation}
The notation \(U_X,D_X\) is as in \eqref{eq:cubic-small-even-connection};
put \(r=(3q)^{1/2}\).  This calculation splits the two nonzero eigenlines
as well as the zero-primary block and checks the rational calculation
independently.

\begin{proposition}[Small-even cubic block reduction]
\label{prop:cubic-block-data}
\coverage{fragment}%
\lean{cubicSmallEven_blockReduction,%
  cubicZeroBlock_modifiedResidue_indicialPolynomial,%
  cubicSmallEven_normalizedGauge_exists_and_unique,%
  cubicSmallEven_normalizedGauge_coefficients,%
  cubicExponentDifference_nonintegral}%
\imports{Beauville:quantum-products}%
The characteristic and minimal polynomials of \(U_X\) are both
\(T^2(T^2-108q)\).  After adjoining \(r\), its eigenvalues are
\(6r,-6r,0,0\), and the zero generalized eigenspace is one rank-two Jordan
block.  A normalized gauge regular in \(z\) gives
\begin{equation}
 \begin{aligned}
 z^2\partial_z\widetilde y_+&=(6r+O(z^2))\widetilde y_+,\\
 z^2\partial_z\widetilde y_-&=(-6r+O(z^2))\widetilde y_-,\\
 z^2\partial_z
 \begin{pmatrix}\widetilde y_3\\\widetilde y_4\end{pmatrix}
 &=\bigl[N+zA_0+z^2A_1+O(z^3)\bigr]
 \begin{pmatrix}\widetilde y_3\\\widetilde y_4\end{pmatrix},
 \end{aligned}
 \label{eq:cubic-integral-block-reduction}
\end{equation}
where
\begin{equation}
 N=\begin{pmatrix}0&2\\0&0\end{pmatrix},\quad
 A_0=\begin{pmatrix}-19/18&0\\0&19/18\end{pmatrix},\quad
 A_1=\begin{pmatrix}0&-14/(81r^2)\\-8/81&0\end{pmatrix}.
 \label{eq:cubic-shared-block-data}
\end{equation}
\end{proposition}

\begin{proof}
The matrix
\[
 C=\begin{pmatrix}
 6r^3&-6r^3&0&-7r^2\\7r^2&7r^2&-2r^2&0\\
 3r&-3r&0&1\\1&1&1&0
 \end{pmatrix},
 \qquad \det C=-486r^5,
\]
puts \(U_X\) into
\[
 C^{-1}U_XC=
 \begin{pmatrix}6r&0&0&0\\0&-6r&0&0\\0&0&0&2\\0&0&0&0\end{pmatrix}.
\]
Partition as \(1\mid1\mid2\).  At each positive order in \(z\), the
off-diagonal gauge coefficient is the unique solution of a Sylvester
equation between blocks with disjoint scalar spectra.  The gauge and inverse
are regular in \(z\); explicitly, it has the normalized form
\begin{equation}
 G(z)=I+\sum_{n\ge1}G_nz^n,
 \qquad G_n\text{ block-off-diagonal}.
 \label{eq:cubic-integral-block-gauge}
\end{equation}
To extract the coefficient used in the residue, write
\(y=CG(z)\widetilde y\), \(J=C^{-1}U_XC\) and \(B=C^{-1}D_XC\).
Label the blocks by \(+,-,0\).  If the transformed coefficient is
\(J+z\widehat A_0+z^2\widehat A_1+\cdots\), then
\[
 (J+zB)G-z^2G'=G(J+z\widehat A_0+z^2\widehat A_1+\cdots).
\]
Thus \((\widehat A_0)_{ii}=B_{ii}\), while the first Sylvester equations are
\[
 J_i(G_1)_{ij}-(G_1)_{ij}J_j=-B_{ij}\qquad(i\ne j).
\]
At the next order,
\[
 \widehat A_1=[J,G_2]+BG_1-G_1-G_1\widehat A_0.
\]
Since \(G_1,G_2\) are block off-diagonal and \(\widehat A_0\) is block
diagonal, the zero block is simply
\[
 A_1=(\widehat A_1)_{00}
 =B_{0+}(G_1)_{+0}+B_{0-}(G_1)_{-0}
 =\begin{pmatrix}0&-14/(81r^2)\\-8/81&0\end{pmatrix}.
\]
This computes the needed coefficient without solving for \(G_2\), and proves
\eqref{eq:cubic-shared-block-data}.
\end{proof}

